\setcounter{paragraph}{0}
\paragraph{Thay đổi tần số sóng}
\Anlythuyet{
\circEX{1}\ \textbf{Hai đầu cố định hoặc hai đầu tự do} 
\begin{align*}
\ell=k\dfrac{\lambda}{2}=k\dfrac{v}{2f}\Rightarrow \boder{f_k=k.\dfrac{v}{2\ell}} 
\end{align*}
\begin{itemize}
\item Đặt: $f_0=\dfrac{v}{2\ell}$ là tần số nhỏ nhất. Vậy $\boder{f_k=kf_0}$ ($k=1,\ 2,\ 3...$)\\
\lefthand\ Tần số để tạo ra sóng dừng trên sợi dây hai đầu cố định bằng \textit{số nguyên} lần tần số nhỏ nhất.
\item Gọi $f_k$ và $f_{k+1}$ là hai tần số liên tiếp cho sóng dừng trên dây hai đầu cố định thì  \begin{align*}\boder{f_0=f_{k+1}-f_k}\end{align*}
\end{itemize}
\circEX{2}\ \textbf{Một đầu cố định, một đầu tự do} 
\begin{align*}
\ell=(2k+1)\dfrac{\lambda}{4}=(2k+1)\dfrac{v}{4f}\Rightarrow \boder{f_k=(2k+1)\dfrac{v}{4\ell}}
\end{align*}
\begin{itemize}
\item Đặt: $f_0=\dfrac{v}{4\ell}$ gọi là tần số nhỏ nhất. Vậy $\boder{f=(2k+1)f_0}$  ($k=0,\ 1,\ 2,\ 3...$) \\
\lefthand\ Tần số để tạo ra sóng dừng trên sợi dây một đầu cố định, một đầu tự do bằng \textit{số lẻ} lần tần số nhỏ nhất.
\item Gọi $f_k$ và $f_{k+1}$ là hai tần số liên tiếp cho sóng dừng trên dây một đầu cố định, một đầu tự do thì \begin{align*}\boder{f_0=\dfrac{f_{k+1}-f_k}{2}}\end{align*}
\end{itemize}
\Note{
Nếu cho biết $f_1\le f\le f_2$ thì thay biểu thức f vào điều kiện giới hạn nói trên để tìm ra giá trị k.
}
}
\Anlythuyet{
\paragraph{Phương trình sóng tại một điểm M cách điểm phản xạ cố định một đoạn d}
\immini{\begin{itemize}
\item Phương trình sóng tới và sóng phản xạ tại B:
\begin{align*}
u_B=A\cos\omega t;\ u_B'=-A\cos\omega t=A\cos(\omega t-\pi)
\end{align*}
\item Phương trình sóng tới và sóng phản xạ tại M cách B một khoảng d:
\begin{align*}
u_M&=A\cos(\omega t+\dfrac{2\pi d}{\lambda});\\  
u'_M&=A\cos(\omega t-\dfrac{2\pi d}{\lambda}-\pi)
\end{align*}
\end{itemize}}{\begin{tikzpicture}[scale=1,thick,>=stealth']
\fill [pattern = north east lines] (8,-0.4) rectangle (8.2,0.4);
\fill [pattern = north east lines] (8,-2.2) rectangle (8.2,-1.3);
\draw (1,0.5) cos (2,0) sin (3,-0.5) cos (4,0) sin(5,0.5) cos(6,0) sin(7,-0.5) cos(8,0);
\draw[line width=0.4pt] (1,0.4) cos (2,0) sin (3,-0.4) cos (4,0) sin(5,0.4) cos(6,0) sin(7,-0.4) cos(8,0);
\draw[line width=0.4pt] (1,0.3) cos (2,0) sin (3,-0.3) cos (4,0) sin(5,0.3) cos(6,0) sin(7,-0.3) cos(8,0);
\draw[line width=0.4pt] (1,0.2) cos (2,0) sin (3,-0.2) cos (4,0) sin(5,0.2) cos(6,0) sin(7,-0.2) cos(8,0);
\draw[line width=0.4pt] (1,0.1) cos (2,0) sin (3,-0.1) cos (4,0) sin(5,0.1) cos(6,0) sin(7,-0.1) cos(8,0);
\draw[dashed,line width=1pt] (1,-0.5) cos (2,0) sin (3,0.5) cos (4,0) sin(5,-0.5) cos(6,0) sin(7,0.5) cos(8,0) ;
\draw[dashed,line width=0.4pt] (1,-0.4) cos (2,0) sin (3,0.4) cos (4,0) sin(5,-0.4) cos(6,0) sin(7,0.4) cos(8,0) ;
\draw[dashed,line width=0.4pt] (1,-0.3) cos (2,0) sin (3,0.3) cos (4,0) sin(5,-0.3) cos(6,0) sin(7,0.3) cos(8,0) ;
\draw[dashed,line width=0.4pt] (1,-0.2) cos (2,0) sin (3,0.2) cos (4,0) sin(5,-0.2) cos(6,0) sin(7,0.2) cos(8,0) ;
\draw[dashed,line width=0.4pt] (1,-0.1) cos (2,0) sin (3,0.1) cos (4,0) sin(5,-0.1) cos(6,0) sin(7,0.1) cos(8,0) ;
\fill[black] (2,0) circle (1pt); \fill[black] (4,0) circle (1pt);\fill[black] (6,0) circle (1pt);\fill[black] (8,0) circle (1pt);
\draw[dashed] (1,0)--(8,0);
\draw[dashed,line width=0.5pt](4,1)--(4,0);\draw[dashed,line width=0.5pt](6,1)--(6,0);\draw[dashed,line width=0.5pt](7,1)--(7,0);\draw[<->](4,0.7)--(6,0.7);\draw[<->](6,0.7)--(7,0.7);\node [above] at (5,0.7) {$\dfrac{\lambda}{2}$}; \node [above] at (6.5,0.7) {$\dfrac{\lambda}{4}$};
\draw[line width=1.5pt](1,-1.8)--(8,-1.8); \fill[black] (1,-1.8) circle (1pt); \fill[black] (4.5,-1.8) circle (1pt); \draw[dashed] (4.5,-1.8)--(4.5,-1.1); \draw[<->](4.5,-1.6)--(8,-1.6);\node [above] at (6.25,-1.6) {d}; \node [below] at (1,-1.8) {A}; \node [below] at (4.5,-1.8) {M}; \node [right] at (8.15,-1.8) {B};
\draw[line width=1.5pt](8,-0.4)--(8,0.4);
\draw[line width=1.5pt] (8,-2.2)--(8,-1.3);
\draw[->,line width=1pt](2,0.3)--(2,0);\node [above] at (2,0.3) {\normalsize{nút sóng}};
\draw[->,line width=1pt](3,0.8)--(3,0.5);\node [above] at (3,0.8) {\normalsize{bụng sóng}};
\draw[->,line width=1pt](4.2,-0.6)--(4.8,-0.2);\node [below] at (4.2,-0.6) {\normalsize{bó sóng}};
\draw[->,black,line width=1pt] (1,-1.4) cos (1.1,-1.5) sin (1.2,-1.6) cos (1.3,-1.5) sin(1.4,-1.4) cos(1.5,-1.5) sin(1.6,-1.6) cos(1.7,-1.5)--(2.2,-1.5);
\node [above] at (2,-1.5) {\normalsize{Sóng tới}};
\draw[->,black,line width=1pt] (8,-2.2) cos (7.9,-2.1) sin (7.8,-2) cos (7.7,-2.1) sin(7.6,-2.2) cos(7.5,-2.1) sin(7.4,-2) cos(7.3,-2.1) sin (7.2,-2.2) cos (7.1,-2.1)--(6.8,-2.1); \node [below] at (6.7,-2.1) {\normalsize{Sóng phản xạ}};
\end{tikzpicture}}
\begin{itemize}
\item Phương trình sóng dừng tại M
\begin{align*}
u=u_M+u_M'=2A\cos(\dfrac{2\pi d}{\lambda}+\dfrac{\pi}{2})\cos(\omega t-\dfrac{\pi}{2})
\end{align*}
\item Biên độ sóng tại M: $\boder{A_M=2A\abs{\cos(\dfrac{2\pi d}{\lambda}+\dfrac{\pi}{2})}=2A\abs{\sin(\dfrac{2\pi d}{\lambda})}}$
\end{itemize}
\paragraph{Phương trình sóng tại một điểm M cách điểm phản tự do một đoạn d}
\begin{itemize}
\item Phương trình sóng tới và sóng phản xạ tại B
\begin{align*}u_B=u'_B=A\cos\omega t\ct{ với }\omega=\dfrac{2\pi}{T}=2\pi f
\end{align*}
\item Phương trình sóng tới và sóng phản xạ tại M cách B một khoảng d
\begin{align*}
u_M=A\cos(\omega t+\dfrac{2\pi d}{\lambda});\ u'_M=A\cos(\omega t-\dfrac{2\pi d}{\lambda})
\end{align*}
\item Phương trình sóng dừng tại M
\begin{align*}
u=u_M+u_M'=2A\cos(\dfrac{2\pi d}{\lambda})\cos\omega t
\end{align*}
\item Biên độ sóng: \bth{\boder{A_M=2A\abs{\cos(\dfrac{2\pi d}{\lambda})}}}
\end{itemize}
\begin{note}
Xác định tốc độ truyền sóng từ phương trình
\begin{align*}
\boder{\ct{Tốc độ truyền sóng}=\dfrac{\ct{hệ số của t}}{\ct{hệ số của x}}}
\end{align*}
\end{note}
}